\chapter{Graph Partitioning}

Throughout this section, let $V$ be a finite type representing the vertices, and $G$ be a simple graph on $V$.

\begin{definition}
    \label{def:boundary}
    \lean{GraphPartitioning.boundary}
    \leanok
    The \textbf{Edge Boundary} of a set $S \subseteq V$, denoted $\partial S$, is the set of edges emanating from the set $S$ to its complement:
    \[
    \partial S = \{e \in E(G) \mid \exists a \in S, \exists b \in S^c, e = \{a, b\}\}
    \]
\end{definition}

\begin{definition}
    \label{def:isoperimetricRatio}
    \lean{GraphPartitioning.isoperimetricRatio}
    \uses{def:boundary}
    \leanok
    The \textbf{isoperimetric ratio} of a set $S \subseteq V$, denoted $\theta(S)$, is the ratio of the size of its edge boundary to its own size:
    \[
    \theta(S) := \frac{|\partial S|}{|S|}
    \]
\end{definition}

\begin{definition}
    \label{def:isoperimetricRatioGraph}
    \lean{GraphPartitioning.isoperimetricRatioGraph}
    \uses{def:isoperimetricRatio}
    \leanok
    The \textbf{isoperimetric ratio of the graph} $G$, denoted $\theta(G)$, is defined as the infimum of $\theta(S)$ over all non-empty subsets $S$ whose size does not exceed half the total number of vertices:
    \[
    \theta(G) := \inf \left\{ \theta(S) \ \middle|\ S \subseteq V, \ 0 < |S| \le \frac{|V|}{2} \right\}
    \]
\end{definition}

\begin{theorem}
    \label{thm:isoperimetric_positive_iff_connected}
    \lean{GraphPartitioning.isoperimetric_positive_iff_connected}
    \uses{def:isoperimetricRatioGraph}
    The isoperimetric ratio of $G$ is strictly positive if and only if $G$ is a connected graph:
    \[
    0 < \theta(G) \iff G.\text{Connected}
    \]
\end{theorem}

\begin{proof}
    Left as an exercise to the Lean compiler. (will write soon)
\end{proof}